\documentclass[11pt,a4paper]{article}

\usepackage[T1]{fontenc}
\usepackage[utf8]{inputenc}
\usepackage[ngerman]{babel}
\usepackage{lmodern}
\usepackage{amsmath,amssymb,amsthm,mathtools}
\usepackage{geometry}
\usepackage{booktabs}
\usepackage{hyperref}
\usepackage{microtype}
\usepackage{enumitem}

\geometry{margin=2.7cm}

\title{Vierquadrate-Zerlegungen von Primzahlen\\
	und zyklische Schließbarkeit im Sehnenmodell}
\author{Thomas Hoffbauer}
\date{\today}

\newtheorem{satz}{Satz}
\newtheorem{definition}{Definition}
\newtheorem{lemma}{Lemma}
\newtheorem{korollar}{Korollar}
\theoremstyle{remark}
\newtheorem{bemerkung}{Bemerkung}
\newtheorem{beispiel}{Beispiel}

\begin{document}
	\maketitle
	
	\begin{abstract}
		Wir untersuchen Vierquadrate-Zerlegungen von Primzahlen in einem zyklischen Sehnenmodell. 
		Ausgehend vom Vier-Quadrate-Satz von Lagrange betrachten wir jede Darstellung
		\[
		p = E^2 + A^2 + B^2 + C^2
		\]
		als geordnete oder ungeordnete Viererstruktur und fragen, ob die zugehörigen Längen als Sehnen eines Kreises so realisiert werden können, dass sich die vier zugehörigen Mittelpunktswinkel zu \(2\pi\) schließen. 
		Dies führt zu einer zusätzlichen geometrischen Verträglichkeitsbedingung innerhalb der arithmetisch garantierten Vierquadrate-Darstellungen. 
		Numerisch zeigt sich für die ersten 1000 Primzahlen, dass die einrastfähigen Darstellungen im Mittel deutlich entropiereicher, weniger konzentriert und balancierter sind als die nicht-einrastfähigen. 
		Ferner wird ein Vergleich zwischen einer linearen und einer quadratischen Gewichtung durchgeführt. Dieser zeigt, dass die quadratische Gewichtung die Schließungsbedingung spürbar verschärft.
	\end{abstract}
	
	\section{Einleitung}
	
	Der Vier-Quadrate-Satz von Lagrange besagt, dass jede natürliche Zahl als Summe von vier Quadraten darstellbar ist. 
	Damit besitzt insbesondere jede Primzahl \(p\) mindestens eine Darstellung der Form
	\[
	p = E^2 + A^2 + B^2 + C^2.
	\]
	Diese klassische Existenz ist rein arithmetisch. 
	Die vorliegende Untersuchung fragt nach einer zusätzlichen geometrischen Struktur solcher Darstellungen.
	
	Die Grundidee besteht darin, die vier Beiträge \(E,A,B,C\) nicht nur als additive Daten zu lesen, sondern sie in ein Kreis- beziehungsweise Sehnenmodell einzubetten. 
	Dabei entsteht zu jeder Vierquadrate-Darstellung eine zyklische Konfiguration, deren Schließbarkeit als zusätzliche Selektionsbedingung interpretiert werden kann. 
	So wird die Menge aller Vierquadrate-Zerlegungen einer Primzahl durch ein geometrisches Kriterium verfeinert.
	
	Der konzeptionelle Kern lautet:
	
	\begin{quote}
		Lagrange liefert die arithmetische Universalität; das Sehnenmodell führt eine geometrische Verträglichkeitsbedingung ein.
	\end{quote}
	
	\section{Vierquadrate-Zerlegungen}
	
	\begin{definition}
		Für \(n\in\mathbb{N}\) sei
		\[
		\mathcal{Q}(n)
		=
		\left\{
		(E,A,B,C)\in\mathbb{Z}_{\ge 0}^4
		\;\middle|\;
		E\ge A\ge B\ge C,\quad
		E^2+A^2+B^2+C^2=n
		\right\}
		\]
		die Menge der ungeordneten Vierquadrate-Darstellungen von \(n\).
	\end{definition}
	
	\begin{satz}[Lagrange]
		Für jede natürliche Zahl \(n\) gilt
		\[
		\mathcal{Q}(n)\neq\emptyset.
		\]
	\end{satz}
	
	\begin{bemerkung}
		Im vorliegenden Artikel werden numerisch ungeordnete Repräsentanten
		\[
		E\ge A\ge B\ge C\ge 0
		\]
		verwendet. Für zyklische Fragestellungen kann später auch die geordnete Variante mit allen verschiedenen Permutationen interessant sein.
	\end{bemerkung}
	
	\section{Das zyklische Sehnenmodell}
	
	\subsection{Zwei Gewichtungsregime}
	
	Zu einer Darstellung \((E,A,B,C)\in\mathcal{Q}(n)\) ordnen wir vier Längen \(l_1,l_2,l_3,l_4\) zu. 
	Wir betrachten zwei Varianten:
	
	\begin{enumerate}[label=\arabic*)]
		\item \textbf{Harte Gewichtung}
		\[
		l_1=E^2,\qquad l_2=A^2,\qquad l_3=B^2,\qquad l_4=C^2.
		\]
		
		\item \textbf{Weiche Gewichtung}
		\[
		l_1=E,\qquad l_2=A,\qquad l_3=B,\qquad l_4=C.
		\]
	\end{enumerate}
	
	Die harte Gewichtung verstärkt große Komponenten überproportional, die weiche Gewichtung behandelt sie linear.
	
	\subsection{Schließungsbedingung}
	
	Für eine Sehne der Länge \(l\) in einem Kreis vom Radius \(r\) gilt
	\[
	\theta(l;r)=2\arcsin\!\left(\frac{l}{2r}\right).
	\]
	Eine Vierer-Konfiguration heißt \emph{einrastfähig} oder \emph{fit}, wenn es einen Radius \(r>0\) gibt mit
	\[
	\theta(l_1;r)+\theta(l_2;r)+\theta(l_3;r)+\theta(l_4;r)=2\pi.
	\]
	Äquivalent:
	\[
	\sum_{i=1}^4 \arcsin\!\left(\frac{l_i}{2r}\right)=\pi.
	\]
	
	\begin{definition}
		Existiert ein Radius \(r>0\), der die Gleichung
		\[
		\sum_{i=1}^4 \arcsin\!\left(\frac{l_i}{2r}\right)=\pi
		\]
		erfüllt, so nennen wir ihn die \emph{Einrastzahl} der Darstellung und schreiben
		\[
		\kappa(E,A,B,C)=r.
		\]
	\end{definition}
	
	\begin{lemma}
		Die Funktion
		\[
		f(r)=\sum_{i=1}^4 \arcsin\!\left(\frac{l_i}{2r}\right)
		\]
		ist auf ihrem Definitionsbereich streng monoton fallend.
	\end{lemma}
	
	\begin{proof}
		Jeder Summand
		\[
		r\mapsto \arcsin\!\left(\frac{l_i}{2r}\right)
		\]
		ist für \(r>\frac{l_i}{2}\) streng fallend, weil \(r\mapsto \frac{l_i}{2r}\) streng fällt und \(\arcsin\) auf \((0,1)\) streng wächst. Die Summe streng fallender Funktionen ist wieder streng fallend.
	\end{proof}
	
	\begin{korollar}
		Falls \(\kappa(E,A,B,C)\) existiert, dann ist sie eindeutig.
	\end{korollar}
	
	\section{Kritischer Defekt}
	
	\begin{definition}
		Sei
		\[
		r_{\min}=\frac12\max(l_1,l_2,l_3,l_4).
		\]
		Dann definieren wir den \emph{kritischen Defekt} durch
		\[
		\delta(E,A,B,C)=\pi-f(r_{\min}),
		\qquad
		f(r)=\sum_{i=1}^4 \arcsin\!\left(\frac{l_i}{2r}\right).
		\]
	\end{definition}
	
	Da \(f(r)\) streng fällt, ist \(f(r_{\min})\) der maximal erreichbare Wert. Also gilt:
	
	\[
	\delta<0 \;\Longrightarrow\; \text{fit},
	\qquad
	\delta=0 \;\Longrightarrow\; \text{Grenzfall},
	\qquad
	\delta>0 \;\Longrightarrow\; \text{non-fit}.
	\]
	
	In der numerischen Praxis wurde fit äquivalent durch \(f(r_{\min})\ge\pi\) mit Toleranz geprüft; fit-Fälle werden dann mit Defekt \(0\) ausgewiesen.
	
	\section{Verteilungsmaße}
	
	Zur Beschreibung der inneren Lastverteilung einer Darstellung verwenden wir die normierten Gewichte
	\[
	p_i=\frac{l_i}{l_1+l_2+l_3+l_4}.
	\]
	
	Daraus definieren wir:
	
	\begin{definition}
		Die \emph{Entropie} einer Darstellung sei
		\[
		H=-\sum_{i=1}^4 p_i\log p_i.
		\]
	\end{definition}
	
	\begin{definition}
		Die \emph{Konzentration} sei
		\[
		K=\sum_{i=1}^4 p_i^2.
		\]
	\end{definition}
	
	\begin{definition}
		Der \emph{Balanceindex} sei
		\[
		\operatorname{bal}
		=
		\frac{\min(l_1,l_2,l_3,l_4)}{\max(l_1,l_2,l_3,l_4)}.
		\]
	\end{definition}
	
	Hohe Entropie und hoher Balanceindex stehen für eine relativ gleichmäßige Lastverteilung, hohe Konzentration für dominante Einzelbeiträge.
	
	\section{Numerische Untersuchung der ersten 1000 Primzahlen}
	
	Es wurden die ersten 1000 Primzahlen
	\[
	2,3,5,\dots,7919
	\]
	analysiert.
	
	\subsection{Harte Gewichtung \(l_i=E_i^2\)}
	
	Im harten Modell besitzen 983 der ersten 1000 Primzahlen mindestens eine fit-Darstellung; 17 Primzahlen besitzen keine fit-Darstellung. 
	Diese 17 Primzahlen sind:
	\[
	5,\;7,\;11,\;17,\;29,\;47,\;53,\;61,\;71,\;83,\;131,\;149,\;197,\;233,\;239,\;257,\;269.
	\]
	
	\subsection{Statistik über alle Darstellungen}
	
	Die statistische Auswertung über sämtliche untersuchten Vierquadrate-Darstellungen der ersten 1000 Primzahlen ergibt im harten Modell:
	
	\begin{center}
		\begin{tabular}{lcccc}
			\toprule
			Gruppe & Anzahl & mittlere Entropie \(H\) & mittlere Konzentration \(K\) & mittlere Balance \\
			\midrule
			fit     & 11046 & 1.1898 & 0.3262 & 0.1977 \\
			non-fit & 77685 & 0.8324 & 0.5302 & 0.0454 \\
			\bottomrule
		\end{tabular}
	\end{center}
	
	Damit zeigen fit-Darstellungen im Mittel:
	
	\begin{itemize}
		\item deutlich \emph{höhere Entropie},
		\item deutlich \emph{geringere Konzentration},
		\item deutlich \emph{größere Balance}.
	\end{itemize}
	
	Dies ist einer der zentralen numerischen Befunde des Modells.
	
	\subsection{Interpretation}
	
	Die zyklische Schließbarkeit ist damit klar \emph{nicht} bloß eine Umformulierung des Vier-Quadrate-Satzes. 
	Vielmehr wirkt das Sehnenmodell als geometrischer Filter innerhalb der von Lagrange gelieferten Darstellungen. 
	Es bevorzugt solche Vierquadrate-Zerlegungen, bei denen die Last nicht zu stark auf eine dominante Hauptkomponente konzentriert ist.
	
	\section{Fast einrastende non-fit-Fälle}
	
	Unter den 17 non-fit-Primzahlen im harten Modell gibt es sehr unterschiedliche Defektgrößen. 
	Einige Fälle sind deutlich non-fit, andere liegen extrem nahe an der Schließung.
	
	Besonders kleine Defekte wurden etwa für folgende Primzahlen gefunden:
	\[
	197,\qquad 61,\qquad 239,\qquad 233,\qquad 149.
	\]
	
	Ein besonders markanter Grenzfall ist \(p=197\) mit der Darstellung
	\[
	(9,8,6,4),
	\]
	für die numerisch
	\[
	\delta \approx 3.93\times 10^{-4}
	\]
	gefunden wurde. 
	Dies zeigt, dass die Klasse der non-fit-Darstellungen selbst eine feine innere Struktur besitzt: 
	Sie reicht von stark rigiden bis zu nahezu schließbaren Konfigurationen.
	
	\section{Vergleich zwischen harter und weicher Gewichtung}
	
	Im weichen Modell
	\[
	l_i=E_i
	\]
	besitzen 993 der ersten 1000 Primzahlen mindestens eine fit-Darstellung. 
	Der Vergleich ergibt somit:
	
	\begin{center}
		\begin{tabular}{lc}
			\toprule
			Größe & Wert \\
			\midrule
			fit im harten Modell & 983 \\
			fit im weichen Modell & 993 \\
			Verbesserungen hard \(\to\) soft & 10 \\
			Verschlechterungen hard \(\to\) soft & 0 \\
			in beiden Modellen non-fit & 7 \\
			\bottomrule
		\end{tabular}
	\end{center}
	
	Dies zeigt eindeutig:
	
	\begin{quote}
		Die quadratische Gewichtung verschärft die geometrische Schließungsbedingung.
	\end{quote}
	
	Große Komponenten werden im harten Modell überproportional verstärkt; dadurch können Darstellungen, die unter linearer Gewichtung noch kompatibel sind, ihre zyklische Schließbarkeit verlieren.
	
	\section{Beste Darstellung pro Primzahl}
	
	Für jede Primzahl wurde zusätzlich eine \emph{beste Darstellung} ausgewählt:
	
	\begin{itemize}
		\item falls fit-Darstellungen existieren: die fit-Darstellung mit kleinster Einrastzahl \(\kappa\),
		\item sonst: die non-fit-Darstellung mit kleinstem Defekt \(\delta\).
	\end{itemize}
	
	Diese Auswahl liefert für jede Primzahl einen kanonischen Repräsentanten ihrer inneren Geometrie im jeweiligen Modell. 
	Im harten Modell zeigt sich dabei häufig, dass die besten fit-Darstellungen eine bemerkenswert gleichmäßige Struktur haben, etwa:
	\[
	(3,3,3,2),\quad
	(4,3,3,3),\quad
	(5,4,4,4),\quad
	(7,6,6,6).
	\]
	Dies stützt erneut die Vermutung, dass die Schließbarkeit eng mit einer balancierten Lastverteilung zusammenhängt.
	
	\section{Schlussfolgerung}
	
	Die vorliegende Untersuchung zeigt, dass Vierquadrate-Darstellungen von Primzahlen durch ein zyklisches Sehnenmodell geometrisch verfeinert werden können. 
	Die durch Lagrange garantierte Existenz von Vierquadrate-Zerlegungen bildet dabei die arithmetische Basis; die zusätzliche Kreis-Schließungsbedingung führt eine nichttriviale geometrische Selektivität ein.
	
	Die numerischen Resultate legen nahe:
	
	\begin{enumerate}[label=\arabic*)]
		\item Einrastfähigkeit ist stark mit hoher Entropie und Balance sowie geringer Konzentration korreliert.
		\item Die harte Gewichtung \(l_i=E_i^2\) wirkt als Verstärker struktureller Ungleichverteilung.
		\item Nicht-einrastende Primzahlen bilden keine homogene Klasse, sondern reichen von stark rigiden bis zu nahezu kritischen Grenzfällen.
		\item Die Wahl des Gewichtungsregimes verändert nicht nur die Fit-Quote, sondern teilweise auch die optimale Darstellung einer Primzahl.
	\end{enumerate}
	
	Damit entsteht aus der klassischen Vierquadrate-Theorie eine arithmetisch-geometrische Verfeinerung, in der nicht nur die Existenz von Darstellungen zählt, sondern auch deren innere zyklische Verträglichkeit.
	
	\section*{Ausblick}
	
	Naheliegende Fortsetzungen dieser Arbeit sind:
	
	\begin{itemize}
		\item Untersuchung geordneter statt nur ungeordneter Vierquadrate-Tupel,
		\item Vergleich der Klassen \(p\equiv 1 \pmod 4\) und \(p\equiv 3 \pmod 4\),
		\item asymptotische Analyse der Fit-Quote für größere Primzahlbereiche,
		\item Untersuchung anderer Gewichtungsregime \(l_i=E_i^\alpha\),
		\item quaternionische Interpretation der Viererstruktur über
		\[
		N(E+Ai+Bj+Ck)=E^2+A^2+B^2+C^2.
		\]
	\end{itemize}
	
\end{document}